\chapter{Time-varying dynamic algorithm review}

\section{Multi-subjects Time-varying graphical Lasso}

A network is composed of a group of nodes connected to each other by edges. 
Whether two nodes are connected or not depends on their dependence conditional on other nodes.
This is straightforward because an edge only reflects the relationship between two nodes regardless of the states of the rest. 
For Gaussian observations, a direct statistic showing this conditional dependence is the precision matrix, i.e. the inverse covariance matrices, 
whereby two nodes are conditionally independent to each other if and only if their corresponding precision matrix entries are zero\footnote{See Subsection \ref{subsec: pcorrandprec}}.
Based on this property, 
Hallac et al.\cite{hallac2017networkinferencetimevaryinggraphical} proposed Time-varying graphical lasso (TVGL).  
This is a method fully leveraging the time series observations i.i.d. from $\mathcal{N}(0,\Sigma(t))$ at different time stamps $t_1,\cdots,t_n$  
to acquire penalized likelihood estimator of precision matrices $\Theta=(\Theta_1, \cdots, \Theta_n)$
and hence to infer the snapshots of a time-varying network.
The method establishes the minimization goal by constructing the loss function consists of a negative log-likelihood with 
the $l_1$ spatial penalty on $\Theta$ and a convex spatial correlation penalty function on the difference between two adjacent precision matrices $\Theta_i$, $\Theta_{i+1}$. 
In addition, Hallac et al.\cite{hallac2017networkinferencetimevaryinggraphical} uses Alternating Direction Method of
Multipliers (ADMM) to derive an efficient algorithm to solve the optimization problem. 

TVGL infers the time-varying network from raw observations, 
whereas it is possible in real-life cases, for each time stamp, 
only sample covariance matrices respectively from different subjects are observable and the sample sizes of observations from each subject are unknown.
In abstract research backgrounds, we may deem a time series of positive definite real matrices as snapshots of a time-varying covariance matrix, 
with possibly multiple matrices for each time stamp. 
The typical research questions include inferring time varying cell communication network from cell communication probability matrix, 
and discovering aging dynamic of brain structure connection. 
In this way, the observations consist of multiple sample covariance matrices but not single one per time stamp, 
on which TVGL can't perform directly.
In order to make TVGL applicable to the aforementioned cases, 
we adapt TVGL to a multi-subjects version called msTVGL and give the corresponding ADMM solver for $l_1$ and $l_2$ spatial penalty on the basis of the work of Hallac et al.. 
The adaption is equivalent to using the average of all observed covariance matrices for each time stamp. 
The msTVGL method is also applicative when in each time stamp the samples are collected from different subjects 
and can't be mix up directly as a whole considering heterogeneity between subjects. 
For example, when finding the aging dynamic of gene interaction, 
if two independent genes happen to express proactively or inactively at the same time in a subject, 
and cells from this subject are much more than the cells from the other subjects, 
the mixed data will contribute to a false high correlation between the two genes, 
thereby allowing TVGL to infer the existence of an unnecessary network edge.
However, the averaged within-subject covariance matrices used by msTVGL will not magnify the heterogeneity between subjects with unequal sample size, 
and hence reduce the strength of unreal network edges. 
Furthermore, the prominent difference that msTVGL distinguishes from TVGL is that 
msTVGL gives the partial correlation matrix at each time stamp to show the connection strength of nodes, 
and is able to show the change from $t_i$ to $t_{i+1}$
using the difference between the two partial correlation matrices. 

% Current TVGL accepts one sample covariance matrix for each time stamp. 
% It suits cases where the target of interest is a single subject. 
% However, in real-life situations, usually it is difficult to conduct continuous follow-up on the same subject. 
% We instead focus on multiple subjects with the identical time-varying pattern to acquire a relatively complete time series observation set. 
% Researches in large-scale epidemiological surveys, multi-center clinical trials or environmental health monitoring 
% are vivid examples. 
% Of note we can still approximately get the sample covariance matrix of the untractable entity by 
% i) mixing up observations from subjects as a whole for each time stamp; 
% or ii) adding up respective sample covariance matrices with weights $\frac{n_{ij}}{n_i}$, $1\leq j\leq k$, 
% if there are $k$ subjects and $n_i=\sum_{j=1}^{k}n_{ij}$ samples in total for time stamp $t_i$.  
% It worth mentioning that this will most effectively incorporate the noise into the results, 
% including imbalanced sample size for each subject and different methods to collect observations, 
% potentially causing the outcomes to deviate from the actual situation 
% and hence preventing TVGL from accurate inference of the time-varying network.
% Therefore, in order to make TVGL applicable to the aforementioned cases, 
% we adapt TVGL to a multi-subjects version called msTVGL and give the corresponding ADMM solver on the basis of the work of Hallac et al..

Below we first introduce the preliminary on the precision matrix and partial correlation, 
and then derive the updating algorithm of msTVGL with ADMM.

\subsection{Precision matrix and partial correlation}
\label{subsec: pcorrandprec}
In this part, we introduce the definition of precision matrix and partial correlation. 
Some useful properties of them are discussed. 
Below we suppose random variables $X=(X_1,X_2,\cdots,X_n)$ are jointly Gaussian with 0 mean and covariance $\Sigma$,  
and denote $X-\{X_i\}$ by $V_i$. 

\begin{definition}(\textbf{Precision matrix})

    The precision matrix $\Theta$ of multivariate Gaussian random variables $(X_1,\cdots,X_n)$ is the inversion of the population covariance $\Sigma$.
\end{definition}

The precision matrix of multivariate Gaussian random variables can be used to test the conditional independence. 
In fact, given values of all other identities, 
any two variables $X_i$ and $X_j$ are conditionally independent if $\Theta_{ij}=0$. 
This is shown in the following theorem.
\begin{theorem}
    \label{thm: ch1_conind}
    $X_i$ and $X_j$ are conditionally independent given $V_{i,j}$ if and only if $\Theta_{ij}=\Theta_{ji}=0$.
\end{theorem}
\begin{proof}
    Without loss of generality we assume $i=1$, $j=2$. For short, we simplify $V_{i,j}$ by $V$.
    The joint Gaussian distribution of $(X_1, X_2)$ condition on $V$ is given by (ignoring the constant) 
    \begin{align*}
        p(X_1,X_2|V)=&|\Theta|^{\frac{1}{2}}\exp\{-\frac{1}{2}[X_1 X_2]\Theta\begin{bmatrix}
            X_1 \\
            X_2 
        \end{bmatrix}\}\\
        =& \text{det}\bigg(\begin{bmatrix}
            \Theta_{11} & \Theta_{12} \\
            \Theta_{21} & \Theta_{22} 
        \end{bmatrix}\bigg)^{\frac{1}{2}}\exp\{-\frac{1}{2}[X_1 X_2]\begin{bmatrix}
            \Theta_{11} & \Theta_{12} \\
            \Theta_{21} & \Theta_{22}
        \end{bmatrix}\begin{bmatrix}
            X_1 \\
            X_2
        \end{bmatrix}\} \\
        =& |\Theta_{11}|^{\frac{1}{2}}\exp\{-\frac{1}{2}X_1\Theta_{11}X_1\}\cdot|\Theta_{22}|^{\frac{1}{2}}\exp\{-\frac{1}{2}X_2\Theta_{22}X_2\} \\
        =& p(X_1|V)\cdotp(X_2|V),
    \end{align*}
    if and only if $\Theta_{12}=\Theta_{21}=0$. Since any $X_i$ and $X_j$ can be moved to the place of $X_1$ and $X_2$ by permutation, this completes the proof. 
\end{proof}

The precision matrix determines the conditional dependence between two Gaussian random variables. 
And the partial correlation measures the relation strength of the two random variables when holding other random variables constant. 
The relation strength is featured by the residuals of $X_i$ and $X_j$ regressing on the rest random variables respectively. 

\begin{definition}(\textbf{Partial correlation})
    
    The partial correlation $\rho_{i,j|V_{i,j}}$ between $X_i$ and $X_j$ given the rest of variables
    is the correlation between the residual resulting from regressing $X_i$ on $V_{i,j}$ and $X_j$ on $V_{i,j}$. 
    More formally, denote residuals by $R_i=X_i-\beta_i^TV_{i,j}$ and $R_j=X_j-\beta^T_jV_{i,j}$. 
    Then, 
    \begin{equation}
        \label{eq: ch1_pcorrdef}
        \rho_{i,j|V_{i,j}} = \frac{\text{cov}(R_i,R_j)}{\sqrt{\text{Var}(R_i)}\sqrt{\text{Var}(R_j)}}.
    \end{equation} 
\end{definition}

Since the residual of $X_i$ regressing on $V_{i,j}$ is the orthogonal projection of the linear space spanned by $X$ on the orthogonal complement of the linear space spanned by $V$, 
it excludes all information that can be explained by the linear combination of $V$, 
whereby it contains the direct linear influence of $X_j$ on $X_i$. 
In addition, the partial correlation has close relationship with the precision matrix. 
\begin{proposition}
    \label{prop: ch1_pcorr}
    Let $\Theta$ be the precision matrix of joint distribution of $X$.
    The partial correlation is equivalent to
    \begin{equation*}
        \rho_{i,j|V_{i,j}}=-\frac{\Theta_{ij}}{\sqrt{\Theta_{ii}\Theta_{jj}}}.
    \end{equation*}
\end{proposition}
\begin{proof}
    Without loss of generality we assume $i=1$, $j=2$. For short, we simplify $V_{i,j}$ by $V$. 
    Since $\beta_1$ and $\beta_2$ are least square solutions for problems
    \begin{align*}
        &\arg\min_{\beta\in\mathbb{R}^{n-2}} \mathbb{E}\|X_1-\beta^TV\|_2^2, \\
        &\arg\min_{\beta\in\mathbb{R}^{n-2}} \mathbb{E}\|X_2-\beta^TV\|_2^2,
    \end{align*}
    and have closed-form $\beta_1=\Sigma_{VV}^{-1}\Sigma_{VX_1}$ and $\beta_2=\Sigma_{VV}^{-1}\Sigma_{VX_2}$ for 0 centered random variables, 
    Equation (\ref{eq: ch1_pcorrdef}) can be written by 
    \begin{align}
        \label{eq: ch1_pcorrprf1}
        \rho_{1,2|V} &= \frac{\text{cov}(R_1,R_2)}{\sqrt{\text{Var}(R_1)}\sqrt{\text{Var}(R_2)}} \notag \\
                     &= \frac{\text{cov}(X_1-\Sigma_{VX_1}^T\Sigma_{VV}^{-1}V, X_2-\Sigma_{VX_2}^T\Sigma_{VV}^{-1}V)}{\sqrt{\text{Var}(X_1-\Sigma_{VX_1}^T\Sigma_{VV}^{-1}V)}\sqrt{\text{Var}(X_2-\Sigma_{VX_2}^T\Sigma_{VV}^{-1}V)}}.
    \end{align}
    For nominator, we have 
    \begin{align*}
        \text{cov}(X_1-\Sigma_{VX_1}^T\Sigma_{VV}^{-1}V, X_2-\Sigma_{VX_2}^T\Sigma_{VV}^{-1}V) &= \mathbb{E}\left[(X_1-\Sigma_{VX_1}^T\Sigma_{VV}^{-1}V)(X_2-\Sigma_{VX_2}^T\Sigma_{VV}^{-1}V)\right] & (\mathbb{E}(R_1)=\mathbb{E}(R_2)=0) \\
                                                                                               &= \Sigma_{X_1X_2} - \Sigma_{VX_2}^{T}\Sigma_{VV}^{-1}\Sigma_{VX_1}.
    \end{align*}
    For denominator, we have 
    \begin{align*}
        \text{Var}(X_1-\Sigma_{VX_1}^T\Sigma_{VV}^{-1}V) &= \mathbb{E}\left[(X_1-\Sigma_{VX_1}^T\Sigma_{VV}^{-1}V)(X_1-\Sigma_{VX_1}^T\Sigma_{VV}^{-1}V)\right] \\
                                                         &= \Sigma_{X_1X_1} -\Sigma_{VX_1}^T\Sigma_{VV}^{-1}\Sigma_{VX_1}.
    \end{align*}
    Hence, Equation (\ref{eq: ch1_pcorrprf1}) equals to 
    \begin{equation}
        \label{eq: ch1_pcorrprf2}
        \frac{\Sigma_{X_1X_2} - \Sigma_{VX_2}^{T}\Sigma_{VV}^{-1}\Sigma_{VX_1}}{\sqrt{\Sigma_{X_1X_1} -\Sigma_{VX_1}^T\Sigma_{VV}^{-1}\Sigma_{VX_1}}\sqrt{\Sigma_{X_2X_2} - \Sigma_{VX_2}^{T}\Sigma_{VV}^{-1}\Sigma_{VX_2}}}.
    \end{equation}
    Let $\Sigma=\begin{bmatrix}
        \Sigma_{X_1X_1} & \Sigma_{X_1X_2} & \Sigma_{X_1V} \\
        \Sigma_{X_2X_1} & \Sigma_{X_2X_2} & \Sigma_{X_2V} \\
        \Sigma_{VX_1} & \Sigma_{VX_2} & \Sigma_{VV}
    \end{bmatrix}=\begin{bmatrix}
        \Sigma_{11} & \Sigma_{12} \\
        \Sigma_{21} & \Sigma_{VV}   
    \end{bmatrix}$. By Schur's formula for block-matrix inversion, 
    $\Theta=\Sigma^{-1}=\begin{bmatrix}
        (\Sigma_{11}-\Sigma_{12}\Sigma_{VV}^{-1}\Sigma_{21})^{-1} & -(\Sigma_{11}-\Sigma_{12}\Sigma_{VV}^{-1}\Sigma_{21})^{-1}\Sigma_{12}\Sigma_{VV}^{-1} \\
        -\Sigma_{VV}^{-1}\Sigma_{21}(\Sigma_{11}-\Sigma_{12}\Sigma_{VV}^{-1}\Sigma_{21})^{-1} & \Sigma_{VV}^{-1}+\Sigma_{VV}^{-1}\Sigma_{21}(\Sigma_{11}-\Sigma_{12}\Sigma_{VV}^{-1}\Sigma_{21})^{-1}\Sigma_{12}\Sigma_{VV}^{-1}
    \end{bmatrix}$.

    With the left-upper corner of $\Theta$ which is
    \begin{align*}
        (\Sigma_{11}-\Sigma_{12}\Sigma_{VV}^{-1}\Sigma_{21})^{-1} &= \bigg(\begin{bmatrix}
            \Sigma_{X_1X_1} & \Sigma_{X_1X_2} \\
            \Sigma_{X_2X_1} & \Sigma_{X_2X_2}
        \end{bmatrix}-\begin{bmatrix}
            \Sigma_{X_1V}\Sigma_{VV}^{-1}\Sigma_{VX_1} & \Sigma_{X_1V}\Sigma_{VV}^{-1}\Sigma_{VX_2} \\
            \Sigma_{X_2V}\Sigma_{VV}^{-1}\Sigma_{VX_1} & \Sigma_{X_2V}\Sigma_{VV}^{-1}\Sigma_{VX_2}
        \end{bmatrix}\bigg)^{-1} \\
                &= \frac{1}{\text{det}}\begin{bmatrix}
                    \Sigma_{X_2X_2}-\Sigma_{X_2V}\Sigma_{VV}^{-1}\Sigma_{VX_2} & -\Sigma_{X_2X_1}+\Sigma_{X_2V}\Sigma_{VV}^{-1}\Sigma_{VX_1} \\
                    -\Sigma_{X_1X_2}+\Sigma_{X_1V}\Sigma_{VV}^{-1}\Sigma_{VX_2} & \Sigma_{X_1X_1}-\Sigma_{X_1V}\Sigma_{VV}^{-1}\Sigma_{VX_1}
                \end{bmatrix} \\
                &= \begin{bmatrix}
                    \Theta_{11} & \Theta_{12} \\
                    \Theta_{21} & \Theta_{22}
                \end{bmatrix},
    \end{align*}
    we observe that Equation (\ref{eq: ch1_pcorrprf2}) is just $-\frac{\Theta_{12}}{\sqrt{\Theta_{11}\Theta_{22}}}$.
    Since any $X_i$ and $X_j$ can be moved to the place of $X_1$ and $X_2$ by permutation, this completes the proof. 
\end{proof}

By the Proposition \ref{prop: ch1_pcorr}, a direct corollary of Theorem \ref{thm: ch1_conind} is obtained. 
\begin{corollary}
    $X_i$ and $X_j$ are conditionally independent given $V_{i,j}$ if and only if $\rho_{i,j|V_{i,j}}=\rho_{j,i|V_{i,j}}=0$.
\end{corollary}

\subsection{Method}

Suppose we have a time series in $\mathbb{R}^d$ over a period of discrete times $t_1,\cdots,t_T$, 
where for each time stamp we have $n_i$ observations. 
We assume the time series is observed from a Gaussian multivariate random vector $X\sim\mathcal{N}(0,\Sigma(t))$. 
To infer the dynamic network from the time series, 
we aim at estimating the sparse inverse covariance matrices $\Theta=\{\Theta_1,\cdots,\Theta_T\}$ from observations. 
This is achievable by minimizing the negative log-likelihood function with a temporal constraint and a sparsity penalty.
Hallac et al. define the loss function by
% \begin{equation*}
%     \mathcal{L} = \sum_{i=1}^{T}\left(-l(\Theta_i) + \lambda\|\Theta_i\|_{\text{od}, 1}\right) + 
%     \beta\cdot\left\{\begin{aligned}
%         &\sum_{i=2}^{T}\psi(\Theta_i-\Theta_{i-1}) & \text{Synchronous Observations}\\
%         &\sum_{i=2}^{T}\Delta_i\psi(\frac{\Theta_i-\Theta_{i-1}}{\Delta_i}) & \text{Asynchronous Observations}
%     \end{aligned}\right.,
% \end{equation*}
\begin{equation*}
    \mathcal{L}_{\text{TVGL}} = \sum_{i=1}^{T}\left(-l(\Theta_i) + \lambda\|\Theta_i\|_{\text{od}, 1}\right) + \beta\sum_{i=2}^{T}\psi(\Theta_i-\Theta_{i-1}), 
\end{equation*}
where $\lambda$, $\beta$ are hyperparameters, $\|\cdot\|_{\text{od}, 1}$ is off-diagonal 1-norm, 
and $\psi(\cdot)$ is a convex penalty function, minimized at $\psi(0)$, 
encouraging similarity between adjacent sparse inverse covariance matrices. 

\subsubsection{Loss function of msTVGL}
Denote the sample covariance matrix by $S$. 
For each $t_i$, the negative log-likelihood function is given by 
\begin{align}
    \label{eq: ch1_loss}
    -l(\Theta_i) &= -\log\left((2\pi)^{-\frac{n_id}{2}}|\Theta_i|^{\frac{n_i}{2}}\exp\left\{-\frac{1}{2}\sum_{j=1}^{n_i}x_j^T\Theta_ix_j\right\}\right) \notag \\
               &= -\log((2\pi)^{-\frac{n_id}{2}}) - \frac{n_i}{2}\log(|\Theta_i|) +\frac{n_i}{2}\sum_{j=1}^{n_i}\text{Tr}(\frac{1}{n_i}x_j^T\Theta_ix_j) \notag \\
               &= -\log((2\pi)^{-\frac{n_id}{2}}) - \frac{n_i}{2}\log(|\Theta_i|) +\frac{n_i}{2}\text{Tr}(\frac{1}{n_i}\sum_{j=1}^{n_i}x_jx_j^T\Theta_i) \notag \\
               &= -\log((2\pi)^{-\frac{n_id}{2}}) - \frac{n_i}{2}\log(|\Theta_i|) +\frac{n_i}{2}\text{Tr}(S_i\Theta_i)
\end{align}
Minimizing negative log-likelihood function encourages $\Theta_i=S_i^{-1}$. 

However, if for each time stamp there are multiple subjects contributing total observations, 
the negative log-likelihood function may not accurate anymore. 
Consider that we test on a certain type of cell in blood samples from multiple subjects. 
In this situation, cells from different subjects compose the set of observations. 
It is feasible to use the combined dataset to estimate the precision matrix but the heterogeneity between each subject should be emphasized. 
In addition, if we only observe several positive definite matrices for each time stamp, 
we can't derive both sample size and the overall "sample covariance matrix". 
Hence, we modify the negative log-likelihood part of the loss function (\ref{eq: ch1_loss}) by 
using sample covariance matrices separately from different subjects while keeping a common precision matrix at each time stamp $t_i$,  
as presented as follows. 
\begin{align}
    \label{eq: ch1_loss_modify1}
    -l(\Theta_i) &= -\log((2\pi)^{-\frac{n_id}{2}}) - \frac{n_i}{2}\log(|\Theta_i|) +\frac{n_i}{2}\text{Tr}(S_i\Theta_i) \notag \\
                 &= \sum_{k=1}^{N_i}-\log((2\pi)^{-\frac{n_{ik}d}{2}})-\frac{n_{ik}}{2}\log(|\Theta_i|)+\frac{n_{ik}}{2}\text{Tr}(S_{ik}\Theta_i),
\end{align}
where there are $N_i$ subjects and $\sum_{k=1}^{N_i}n_{ik}=n_i$.

To reduce the influence of imbalanced sample sizes, 
we further normalize (\ref{eq: ch1_loss_modify1}) by removing the sample sizes $n_{i1},\cdots,n_{iN_i}$. 
This produces the final log-likelihood loss suitable for the multi-subjects case (up to a constant and scale): 
\begin{equation}
    \label{eq: ch1_loss_modify1_normalized}
    -l(\Theta_i) = \sum_{k=1}^{N_i}-\log(|\Theta_i|)+\text{Tr}(S_{ik}\Theta_i).
\end{equation}
To approximately optimize the loss (\ref{eq: ch1_loss_modify1_normalized}), 
we use the mean of sample covariances of all subjects $\frac{1}{N_i}\sum_{k=1}^{N_i}S_{ik}$ denoted by $\bar{S_i}$ and produces a new loss 
\begin{align}
    \label{eq: ch1_loss_modify2_normalized}
    -l(\Theta_i) = -N_i\log(|\Theta_i|)+\text{Tr}(\bar{S_i}\Theta_i).
\end{align}
Of note optimizing (\ref{eq: ch1_loss_modify2_normalized}) produces a precision matrix close to $\bar{S_i}^{-1}$, 
which ignore the difference between subjects.

The spatial penalty function is given by the off-diagonal 1-norm of each precision matrix, 
controlled by a penalty parameter $\lambda$. 
It encourages the sparsity of the precision matrix.

For temporal penalty function $\psi$, 
Hallac et al.\cite{hallac2017networkinferencetimevaryinggraphical} proposed 5 choices, as listed below, 
which are element-wise $l_1$ penalty, group lasso $l_2$ penalty, Laplacian penalty, 
$l_{\infty}$ penalty and row-column overlap penalty. 
Here element-wise $l_1$ penalty $\psi(X)=\sum_{i,j}|X_{i,j}|$ 
encourages the networks at adjacent time points to remain consistent on the majority of the edges, 
while only changing on a few edges. 
Group lasso $l_2$ penalty $\psi(X)=\sum_{j}\|[X]_j\|_2$ 
leads to a complete reconfiguration of the entire network at a few specific time points, 
when maintaining structural stability at the rest time points.
Laplacian penalty $\psi(X)=\sum_{i,j}X_{i,j}^2$ is suitable for smooth time-varying dynamic 
by allowing slight difference between adjacent time stamps and heavily punishing prominent differences. 
$l_{\infty}$ penalty $\psi(X)=\sum_{j}(\max_i|X_{i,j}|)$ is designed for scenarios where change happens only cluster by cluster
because within the largest variation range of all the edges connected to a certain node, 
all edges in this cluster can freely vary without incurring additional penalties.
row-column overlap penalty $\psi(X)=\min_{V:V+V^T=X}\sum_{j}\|[V]_j\|_2$ is suitable for situations of repositioning of a single node to a new set of neighbors at each time stamp.
As we especially focus on research background of time-varying trend of cell-cell communication and brain region interaction, 
we discuss the most biologically meaningful temporal penalty types.

Usually, the time-varying dynamic of cell-cell interaction networks exhibits phased changes. 
Within each phase, the changes are smooth and continuous. 
Between different phases, there are sudden changes in edges and possibly network reorganization.
The reason is that those networks reflect the overall influence from 
the gradual establishment of cell communication patterns during development like immune system maturation during adolescence, 
the slow attenuation of communication intensity or the compensatory enhancement of specific pathways during aging, 
as well as the sudden activation or inhibition of certain signaling axes during disease progression. 
The change of brain structural connectivity network exhibits either a gradual smooth transition or sudden network reconfiguration.
The biological progression includes a continuous but gradually slowing down process of myelination and pruning of fiber bundles (with smooth and continuous changes) from childhood to adulthood;
possible relatively rapid global network reorganization before and after key developmental milestones (such as the maturation of language areas and prefrontal cortex functions).  
In special cases, the brain may suffer from a global connectivity disruptions in the course of neurodegenerative diseases (such as Alzheimer's disease), manifested as phased disintegration of the network.

Based on the aforementioned change characteristics of both types of networks, 
this paper first of all suggests Laplacian penalty functions because it aligns with the smooth and gradual change of the two types of networks, 
especially in studies with densely sampled time points.
Moreover, the element-wise $l_1$ penalty is adopted considering that in both types of networks, 
there are "switch-like" change events between a few key edges. 
For instance, the activation of a specific signaling pathway leads to a significant increase in the communication probability between certain cell types, 
or the sudden myelination of white matter tracts results in enhanced connectivity between certain brain regions.
At last, group lasso $l_2$ penalty works in the study since it is able to sniff out the sudden transition of cell-cell interaction pattern during certain stage in development, 
as well as to catch the large difference between networks when the gap between time points are huge, 
for example in studies across very rough time labels like "neonate", "adolescent", "adult" and "the elder". 
For the rest two penalty types $l_{\infty}$ and row-column overlap, 
since the coordinated changes of edges within a node cluster 
and the behavior of a single node reconnecting all edges at a single time point are extremely rare in cellular interaction networks and brain structural connection networks, 
and may even lack clear biological mechanism explanations, 
we do not intend to adopt them.

The above discussion on the negative loss function, spatial penalty and temporal penalty gives the final loss function of msTVGL,
as shown by (\ref{eq: ch1_loss_modify_overall})
\begin{equation}
    \label{eq: ch1_loss_modify_overall}
    \mathcal{L}_{\text{msTVGL}} = \sum_{i=1}^{T}(-N_i\log(|\Theta_i|)+\text{Tr}(\bar{S_i}\Theta_i) + \lambda\|\Theta_i\|_{\text{od},1}) +
    \beta\sum_{i=2}^{T}\left\{\begin{aligned}
        &\sum_{j,k}|(\Theta_i-\Theta_{i-1})_{j,k}| \\
        &\sum_{j}\|[\Theta_i-\Theta_{i-1}]_j\|_2 \\
        &\sum_{j,k}(\Theta_i-\Theta_{i-1})^2_{j,k}
    \end{aligned}\right..
\end{equation}

\subsubsection{ADMM solution for optimizing task}


\subsubsection{Hyperparameters}
There are two positive hyperparameters controlling penalty strength. 
For sparsity penalty $\lambda\|\Theta\|_{\text{od}, 1}$, 
each off-diagonal element $x$ has gradient\footnote{Subgradient at $x=0$ is set to 0 but not $[-\lambda,\lambda]$.}
\begin{equation*}
    \left\{\begin{aligned}
        &\lambda, x>0 \\
        &0, x=0 \\
        -&\lambda, x<0
    \end{aligned}\right..
\end{equation*}
Since when optimizing convex penalty terms we expect zero gradients, 
the algorithm will push off-diagonal elements to zero hard as $\lambda$ gets large, 
and this result in a very sparse estimated inverse covariance matrix. 
When $\lambda$ is small, we expect $\Theta$ better matches the empirical data, 
but will lead to very dense networks, 
which are less interpretable and often overfit. 

This can be shown mathematically by gradient. 
In fact, when we set temporal penalty to zero. The gradient\footnote{Calculate by Fr\'echet derivative.} of the loss function with respect to $\Theta_t$ is given by 
\begin{align*}
    -\nabla_{\Theta_t}l &= \nabla_{\Theta_t}\left(-\frac{n_t}{2}\log(|\Theta_t|)+\frac{n_t}{2}\text{Tr}(S_t\Theta_t) + \lambda\|\Theta_t\|_{\text{od}, 1}\right) \\
                                  &= -\frac{n_t}{2}\Theta_t^{-1} + \frac{n_t}{2}S_t + (\pm\lambda)_{\text{od}},
\end{align*}
where $(\pm\lambda)_{\text{od}}$ represents a square matrix on $\mathbb{R}^{d\times d}$ with 0 diagonal and $\pm\lambda$ filling its off-diagonal. 
$\nabla_{\Theta_t}-l=0$ is equivalent to $\Theta_t^{-1}=S_t+\frac{2}{n_t}(\pm\lambda)_{\text{od}}$.
We can observe that as $\lambda$ increases, the precision matrix gets sparse from $S_t$. 
In addition, the magnitude of $\lambda$ is associated with sample size at time $t$. 
If sample sizes are extremely imbalanced, 
a common sparsity penalty will cause over penalized precision matrices at time stamps with very few observations. 
This problem is especially often in high dimensional cases.
Hence, to alleviate it, for each time stamp, we usually normalize the log-likelihood function by dividing the sample size. 

For temporal penalty term, $\beta$ controls the transition between adjacent estimated inverse covariance matrices. 
Small $\beta$ will lead to fluctuating $\Theta$ between time stamps 
while large $\beta$ makes transition smooth. 
As $\beta\longrightarrow\infty$, 
the problem becomes a static graphical lasso problem since a constant $\Theta$ is the solution.

\subsubsection{Hyperparameter selection}
To generate an interpretable and biologically meaningful result, 
the algorithm must assign the estimated precision matrices a rational sparsity to highlight significant edges and avoid learning to many noises from observations. 
In addition, the temporal smoothness penalty must be given a proper strength to maintain consecution between adjacent time stamps. 
To achieve this, we introduce the average degree of nodes, Jaccard similarity of adjacent graphs and temporal deviation ratios as standards to restrict the two penalty coefficients. 

The average degree of nodes in a graph is defined by the mean number of edges connected to each node. 
Suppose there are $L$ edges in the graph and $p$ nodes, for an undirected graph, 
the average degree is given by 
\begin{equation*}
    \text{AD} = \frac{2L}{p}, 
\end{equation*}
where the coefficient 2 is due to each edge for a pair is checked only once. 
This measure represents the number of biologically meaningful edges in a gene-gene interaction network. 
Usually, we perceive entries with absolute values larger than $1e-3$ in a precision matrix to be biologically 
meaningful edges. 

The Jaccard similarity index measures the similarity between two sets. 
In the graph connectivity perspective, it represents the overlap rate of the two adjacent edge sets. 
Suppose $E_i$ and $E_{i+1}$ are two edge sets of the two graphs at adjacent time stamps $i$ and $i+1$, 
their Jaccard similarity index is given by 
\begin{equation*}
    J(i,i+1) = \frac{|E_i\bigcap E_{i+1}|}{|E_i\bigcup E_{i+1}|}.
\end{equation*} 
The Jaccard similarity index lies in $[0,1]$. 
When $J(i,i+1)$ is almost 1, the adjacent graphs are almost identical which implies a relatively strong temporal smoothness penalty. 
When $J(i,i+1)$ is close to 0, there are almost no shared edges between the two graphs. 
It can be resulted by a very small temporal smoothness penalty or an extremely strong sparsity penalty, 
or even by both in a complicated biological context. 

The temporal deviation ratio describes the sharp transitions between graphs. 
It is calculated by the average difference of two adjacent precision matrices in Frobenius norm. 
For precision matrices $\Theta_i$ and $\Theta_{i+1}$, the temporal deviation ratio at time shift $i\longrightarrow i+1$ among $T$ time stamps is given by
\begin{equation*}
    \text{TDR}(i, i+1) = \frac{\|\Theta_{i+1}-\Theta_i\|_{\text{Fro}}}{\sum_{t=1}^{T-1}\|\Theta_{t+1}-\Theta_{t}\|_{\text{Fro}}}.
\end{equation*} 
We consider the time shift with a temporal deviation ratio above the $1.25$ times of the average a sharp transition. 

We choose the acceptable hyperparameter combinations by making sure
i) the average degree is between $5\%$ to $15\%$ of node number; 
ii) the average Jaccard similarity is between $0.35$ and $0.8$;
iii) the temporal deviation ratio displays the least but at least 1 sharp transitions among all the attempts. 

The selection procedure is shown in Algorithm \ref{alg: ch1_hyperparametersel}. 

\begin{algorithm}
    \label{alg: ch1_hyperparametersel}
    \caption{Hyperparameter selection with warm start}
    \begin{algorithmic}
        \State 1) $S_{\lambda}\longleftarrow$ Searching space of $\lambda$.
        \State 2) With hyperparameter combination $(\lambda, \beta)$:
        \State \quad 2a) $\text{AD}_{\lambda,\beta}\longleftarrow$ Mean value of average degrees of all precision matrices at all time stamps.
        \State \quad 2b) $\text{TDR}_{\lambda,\beta}\longleftarrow$ Temporal deviation ratios of all precision matrices at all time stamps.
        \State \quad 2c) $\text{ST}_{\lambda,\beta}\longleftarrow$ Sharp transition numbers. 
        \State \quad 2d) $J_{\lambda,\beta}\longleftarrow$ Mean value of Jaccard similarities of all precision matrices at all time stamps.
        \State 3) $\text{AD}_{\text{min}}, \text{AD}_{\text{max}}\longleftarrow$ $5\%, 15\%$ percentiles of node number. 
        \For{$\lambda, i\in \text{Enumerate}(S_{\lambda})$}
            \For{$\beta\in S_{\lambda}[0:i+1]$}
            \State a) Train model with $(\lambda, \beta)$.
            \State b) Compute $\text{AD}_{\lambda,\beta}$, $\text{TDR}_{\lambda,\beta}$, $J_{\lambda,\beta}$
            \EndFor
            \If{$\text{AD}_{\lambda^*,\beta^*}\in[\text{AD}_{\text{min}}, \text{AD}_{\text{max}}]$, $\text{ST}_{\lambda^*,\beta^*}=\arg\min_{\lambda,\beta}\{\text{ST}_{\lambda,\beta}>0\}$, $J_{\lambda^*,\beta^*}\in[0.35,0.8]$}
                \State Accept $(\lambda^*, \beta^*)$
            \EndIf
        \EndFor
        \State \Return All accepted hyperparameters
    \end{algorithmic}
\end{algorithm}